\documentclass[10pt,a4paper]{amsart}
\usepackage{amsmath,amssymb,amsthm,graphicx}
% Retained minimal preamble for the Stacks Project chapter more-algebra.tex (More on Algebra).
% The chapter file pinned-original-source.tex carries only \input{preamble} on its first line, so no
% preamble text lives inside the chapter itself.  The Stacks Project uses one single shared preamble for
% every chapter; the snapshot stages it identically as sources/stacks-chow/preamble.tex and
% sources/stacks-etale-cohomology/preamble.tex.  The block below is copied mechanically, line for line,
% from sources/stacks-chow/preamble.tex lines 14-18 (the Xy-pic setup) and 154-217 (the native theorem
% declarations together with the mathematical macros).  Nothing was added, removed or redefined; the
% wrapper class and the amsmath/amssymb/amsthm/graphicx base packages are supplied by the assembler.
% Chapter: more-algebra.tex -- More on Algebra.
% For commutative diagrams we use Xy-pic
\usepackage[all]{xy}

% We use 2cell for 2-commutative diagrams.
\xyoption{2cell}

% Theorem environments.
%
\theoremstyle{plain}
\newtheorem{theorem}[subsection]{Theorem}
\newtheorem{proposition}[subsection]{Proposition}
\newtheorem{lemma}[subsection]{Lemma}

\theoremstyle{definition}
\newtheorem{definition}[subsection]{Definition}
\newtheorem{example}[subsection]{Example}
\newtheorem{exercise}[subsection]{Exercise}
\newtheorem{situation}[subsection]{Situation}

\theoremstyle{remark}
\newtheorem{remark}[subsection]{Remark}
\newtheorem{remarks}[subsection]{Remarks}

\numberwithin{equation}{subsection}

% Macros
%
\def\lim{\mathop{\mathrm{lim}}\nolimits}
\def\colim{\mathop{\mathrm{colim}}\nolimits}
\def\Spec{\mathop{\mathrm{Spec}}}
\def\Hom{\mathop{\mathrm{Hom}}\nolimits}
\def\Ext{\mathop{\mathrm{Ext}}\nolimits}
\def\SheafHom{\mathop{\mathcal{H}\!\mathit{om}}\nolimits}
\def\SheafExt{\mathop{\mathcal{E}\!\mathit{xt}}\nolimits}
\def\Sch{\mathit{Sch}}
\def\Mor{\mathop{\mathrm{Mor}}\nolimits}
\def\Ob{\mathop{\mathrm{Ob}}\nolimits}
\def\Sh{\mathop{\mathit{Sh}}\nolimits}
\def\NL{\mathop{N\!L}\nolimits}
\def\CH{\mathop{\mathrm{CH}}\nolimits}
\def\proetale{{pro\text{-}\acute{e}tale}}
\def\etale{{\acute{e}tale}}
\def\QCoh{\mathit{QCoh}}
\def\Ker{\mathop{\mathrm{Ker}}}
\def\Im{\mathop{\mathrm{Im}}}
\def\Coker{\mathop{\mathrm{Coker}}}
\def\Coim{\mathop{\mathrm{Coim}}}

% Boxtimes
%
\DeclareMathSymbol{\boxtimes}{\mathbin}{AMSa}{"02}

%
% Macros for moduli stacks/spaces
%
\def\QCohstack{\mathcal{QC}\!\mathit{oh}}
\def\Cohstack{\mathcal{C}\!\mathit{oh}}
\def\Spacesstack{\mathcal{S}\!\mathit{paces}}
\def\Quotfunctor{\mathrm{Quot}}
\def\Hilbfunctor{\mathrm{Hilb}}
\def\Curvesstack{\mathcal{C}\!\mathit{urves}}
\def\Polarizedstack{\mathcal{P}\!\mathit{olarized}}
\def\Complexesstack{\mathcal{C}\!\mathit{omplexes}}
% \Pic is the operator that assigns to X its picard group, usage \Pic(X)
% \Picardstack_{X/B} denotes the Picard stack of X over B
% \Picardfunctor_{X/B} denotes the Picard functor of X over B
\def\Pic{\mathop{\mathrm{Pic}}\nolimits}
\def\Picardstack{\mathcal{P}\!\mathit{ic}}
\def\Picardfunctor{\mathrm{Pic}}
\def\Deformationcategory{\mathcal{D}\!\mathit{ef}}
\usepackage{cleveref}
\crefname{originallemma}{original lemma}{original lemmas}
\Crefname{originallemma}{Original lemma}{Original lemmas}
\makeatletter
\newlabel{more-algebra-lemma-lift-factorization-monic}{{0ALH}{source}{}{sourcepointer-1}{}}
\newlabel{more-algebra-lemma-lift-factorization-monic@cref}{{[originallemma][0][]0ALH}{[1][0][]source}{}{}{}}
\makeatother
\crefname{originallemma}{original lemma}{original lemmas}
\Crefname{originallemma}{Original lemma}{Original lemmas}
\makeatletter
\newlabel{algebra-example-factor-polynomials-etale}{{00UA}{source}{}{sourcepointer-2}{}}
\newlabel{algebra-example-factor-polynomials-etale@cref}{{[originallemma][0][]00UA}{[1][0][]source}{}{}{}}
\makeatother
\crefname{originallemma}{original lemma}{original lemmas}
\Crefname{originallemma}{Original lemma}{Original lemmas}
\makeatletter
\newlabel{lemma-localize-upstairs}{{07M0}{source}{}{sourcepointer-3}{}}
\newlabel{lemma-localize-upstairs@cref}{{[originallemma][0][]07M0}{[1][0][]source}{}{}{}}
\makeatother
\crefname{originallemma}{original lemma}{original lemmas}
\Crefname{originallemma}{Original lemma}{Original lemmas}
\makeatletter
\newlabel{algebra-lemma-map-between-etale}{{00U7}{source}{}{sourcepointer-4}{}}
\newlabel{algebra-lemma-map-between-etale@cref}{{[originallemma][0][]00U7}{[1][0][]source}{}{}{}}
\makeatother
\crefname{originallemma}{original lemma}{original lemmas}
\Crefname{originallemma}{Original lemma}{Original lemmas}
\makeatletter
\newlabel{algebra-lemma-surjective-flat-finitely-presented}{{00U8}{source}{}{sourcepointer-5}{}}
\newlabel{algebra-lemma-surjective-flat-finitely-presented@cref}{{[originallemma][0][]00U8}{[1][0][]source}{}{}{}}
\makeatother
\title{Lifting a factorisation of a monic polynomial along an etale ring map}
\author{The Stacks Project Authors}
\date{stacks-project:more-algebra.tex:local-snapshot-20261012; GNU FDL 1.2 or later}
\begin{document}
\maketitle
\setcounter{subsection}{1}
\noindent\textit{Editorial note.} All delivered mathematical text is the author's own, copied from the pinned Stacks Project chapter file more-algebra.tex; no statement or proof step was rewritten, reordered or removed.\par\medskip
\noindent\textit{Editorial note.} The chapter file pinned-original-source.tex carries only \textbackslash{}input\{preamble\} on its first line, so no preamble text lives inside the chapter. The retained file sources/stacks-more-algebra/retained-preamble.tex is a mechanical, line-for-line copy of the project's single shared preamble as staged in the snapshot (Xy-pic setup plus theorem declarations and mathematical macros). Nothing was added, removed or redefined, and the excerpt serves only as preamble, never as a span coordinate.\par\medskip
\noindent\textit{Editorial note.} Both counted spans start at the \textbackslash{}begin line of their environment, so the statement and the author's proof are delivered whole; surrounding prose of the section is deliberately outside the counted spans and is carried as standard prerequisites instead.\par\medskip
\noindent\textit{Editorial note.} The citation fields (source\_id, source\_version, authors, source\_url, license.name, source\_title) are identical to those of every other card this lane has produced from this chapter, so the batch selector treats them as one source.\par\medskip
\noindent\textit{Allowed standard prerequisites (editorial).}\begin{itemize}
\item Etale ring map and the lifting of factorisations of monic polynomials along it. discussed in the same section, source lines 1740-1842, before the counted statement. Supplies the class of maps along which the factorisation is lifted.
\item Criterion for lifting a factorisation of a monic polynomial modulo an ideal. recalled at the start of the proof, proof line 1854. Provides the statement to which the lemma is reduced.
\item Behaviour of the discriminant under base change, used to keep the factors coprime after the base change. used at proof lines 1860-1890. Supplies the verification that the two factors remain coprime modulo the ideal.
\end{itemize}\medskip
\ifdefined\hypertarget\hypertarget{sourcepointer-1}{}\fi
\noindent\textit{Source reference.} Original lemma 0ALH: focused-seeds/AGP/sources/stacks-snapshot-20261011/tags.txt, source lines 13370--13370. The lemma's own canonical tag, quoted so that readers can resolve the label inside the delivered text; it points at the counted statement itself and carries no proof content.\par\medskip
\ifdefined\hypertarget\hypertarget{sourcepointer-2}{}\fi
\noindent\textit{Source reference.} Original lemma 00UA: focused-seeds/AGP/sources/stacks-snapshot-20261011/tags.txt, source lines 1104--1104. Navigation pointer to a statement the counted proof refers to but does not contain; nothing load-bearing is delegated to it.\par\medskip
\ifdefined\hypertarget\hypertarget{sourcepointer-3}{}\fi
\noindent\textit{Source reference.} Original lemma 07M0: focused-seeds/AGP/sources/stacks-snapshot-20261011/tags.txt, source lines 9710--9710. Navigation pointer to a statement the counted proof refers to but does not contain; nothing load-bearing is delegated to it.\par\medskip
\ifdefined\hypertarget\hypertarget{sourcepointer-4}{}\fi
\noindent\textit{Source reference.} Original lemma 00U7: focused-seeds/AGP/sources/stacks-snapshot-20261011/tags.txt, source lines 1101--1101. Navigation pointer to a statement the counted proof refers to but does not contain; nothing load-bearing is delegated to it.\par\medskip
\ifdefined\hypertarget\hypertarget{sourcepointer-5}{}\fi
\noindent\textit{Source reference.} Original lemma 00U8: focused-seeds/AGP/sources/stacks-snapshot-20261011/tags.txt, source lines 1102--1102. Navigation pointer to a statement the counted proof refers to but does not contain; nothing load-bearing is delegated to it.\par\medskip
\begin{lemma}
\label{lemma-lift-factorization-monic}
Let $A$ be a ring, let $I \subset A$ be an ideal. Let $f \in A[x]$ be a
monic polynomial. Let $\overline{f} = \overline{g} \overline{h}$ be a
factorization of $f$ in $A/I[x]$ such that $\overline{g}$ and $\overline{h}$
are monic and generate the unit ideal in $A/I[x]$. Then there exists an
\'etale ring map $A \to A'$ which induces an isomorphism $A/I \to A'/IA'$
and a factorization $f = g' h'$ in $A'[x]$ with $g'$, $h'$ monic
lifting the given factorization over $A/I$.
\end{lemma}

\begin{proof}
We will deduce this from results on the universal factorization proved
earlier; however, we encourage the reader to find their own proof not
using this trick. Say $\deg(\overline{g}) = n$ and $\deg(\overline{h}) = m$ so
that $\deg(f) = n + m$. Write $f = x^{n + m} + \sum \alpha_i x^{n + m - i}$
for some $\alpha_1, \ldots, \alpha_{n + m} \in A$. Consider the ring map
$$
R = \mathbf{Z}[a_1, \ldots, a_{n + m}]
\longrightarrow
S = \mathbf{Z}[b_1, \ldots, b_n, c_1, \ldots, c_m]
$$
of Algebra, Example \ref{algebra-example-factor-polynomials-etale}.
Let $R \to A$ be the ring map which sends $a_i$ to $\alpha_i$.
Set
$$
B = A \otimes_R S
$$
By construction the image $f_B$ of $f$ in $B[x]$ factors, say
$f_B = g_B h_B$ with $g_B = x^n + \sum (1 \otimes b_i) x^{n - i}$
and similarly for $h_B$.
Write $\overline{g} = x^n + \sum \overline{\beta}_i x^{n - i}$ and
$\overline{h} = x^m + \sum \overline{\gamma}_i x^{m - i}$.
The $A$-algebra map
$$
B \longrightarrow A/I, \quad
1 \otimes b_i \mapsto \overline{\beta}_i, \quad
1 \otimes c_i \mapsto \overline{\gamma}_i
$$
maps $g_B$ and $h_B$ to $\overline{g}$ and $\overline{h}$ in $A/I[x]$.
The displayed map is surjective; denote $J \subset B$ its kernel.
From the discussion in 
Algebra, Example \ref{algebra-example-factor-polynomials-etale}
it is clear that $A \to B$ is etale at all points of $V(J) \subset \Spec(B)$.
Choose $g \in B$ as in Lemma \ref{lemma-localize-upstairs} and
consider the $A$-algebra $B_g$. Since $g$ maps to a unit
in $B/J = A/I$ we obtain also a map $B_g/I B_g \to A/I$ of $A/I$-algebras.
Since $A/I \to B_g/I B_g$ is \'etale, also $B_g/IB_g \to A/I$ is \'etale
(Algebra, Lemma \ref{algebra-lemma-map-between-etale}). Hence there exists an
idempotent $e \in B_g/I B_g$ such that $A/I = (B_g/I B_g)_e$
(Algebra, Lemma \ref{algebra-lemma-surjective-flat-finitely-presented}).
Choose a lift $h \in B_g$ of $e$. Then $A \to A' = (B_g)_h$ with
factorization given by the image of the factorization $f_B = g_B h_B$
in $A'$ is a solution to the problem posed by the lemma.
\end{proof}

\end{document}
